\documentclass[a4paper,10pt]{article}

\usepackage{geometry}
\usepackage[utf8]{inputenc}
\usepackage[T1]{fontenc}
\usepackage{graphicx}
\usepackage{url}\usepackage{multicol}

\geometry{dvips,a4paper,margin=1.5cm}

\title{De Curvis Elasticis}
\author{Leonhard Euler}
\date{1744}

\begin{document}
\setlength{\columnseprule}{1pt}

\maketitle

\section*{Introduction}
\begin{multicols}{2}
Jam pridem summi quique Geometræ agnoverunt, Methodi in hoc Libro traditæ non solum maximum esse usum in ipsa Analysi, sed etiam eam ad resolutionem Problematum physicorum amplissimum subsidium afferre.

Cum enim Mundi universi fabrica sit perfectissima, atque a Creatore sapientissimo absoluta, nihil omnino in mundo contingit, in quo non maximi minimive ratio quæpiam eluceat: quamobrem dubium prorsus est nullum, quin omnes Mundi effectus ex causis finalibus, ope Methodi maximorum \& minimorum æque feliciter determinari queant, atque ex ipsis causis efficientibus.

Hujus rei vero passim tam eximia extant specimina, ut ad veritatis confirmationem pluribus Exemplis omnino non indigeamus; quin potius in hoc erit elaborandum, ut, in quovis Quæstionum naturalium genere, ea investigetur quantitas, quæ maximum minimumve inducat valorem: quod negotium ad Philosophiam potius quam ad Mathesin pertinere videtur.

Cum igitur duplex pateat via effectus Naturæ cognoscendi; altera per causas efficientes, quæ Methodus directa vocari solet; altera causas finales;
Mathematicus utraque pari successu utitur.

Quando scilicet causæ efficientes nimis sunt absconditæ, finales autem nostram
cognitionem minus effugiunt; per Methodum indirectam Quæstio solet resolvi: e contrario autem Methodus directa adhibetur, quoties ex causis efficientibus effectum definire licet.

Inprimis autem opera est adhibenda, ut per utramque viam aditus ad Solutionem aperiatur: sic enim non solum altera Solutio per alteram maxime confirmatur, sed etiam ex utriusque consensu summam percipimus voluptatem.

Hoc modo, curvatura funis seu catenae suspensae duplici via est eruta; altera a priori, ex sollicitationibus gravitatis; altera vero per Methodum maximo-
rum ac minimorum, quoniam funis ejusmodi curvaturam recipere debere intelligebatur, cujus centrum gravitatis infimum obtineret locum. Similiter curvatura radiorum per medium diaphanum varia densitatis transeuntium, tam a priori est determinata, quam etiam ex hoc principio, quod tempore brevissimo
ad datum locum pervenire debeant.

Plurima autem alia similia exempla a Viris Celeberrimis Bernoulliis, aliisque, sunt prolata, quibus tam Methodus solvendi a priori, quam cogni-
tio causarum efficientium maxima accepit incrementa.

\columnbreak
Depuis longtemps déjà, tous les géomètres les plus éminents ont reconnu que la méthode exposée dans ce livre est non seulement d’une très grande utilité en analyse elle-même, mais qu’elle apporte en outre un secours très considérable à la résolution des problèmes physiques.

En effet, puisque la structure de l’univers est tout à fait parfaite et achevée par un créateur très sage, il n’arrive absolument rien dans le monde où n’apparaisse quelque raison de maximum ou de minimum. C’est pourquoi il n’y a aucun doute que tous les effets du monde peuvent être déterminés à partir des causes finales, à l’aide de la méthode des maxima et minima, tout aussi heureusement que par les causes efficientes.

D’ailleurs, on trouve partout des exemples si remarquables de ce fait qu’il n’est nullement nécessaire d’en ajouter de nombreux pour en confirmer la vérité. Il vaut mieux s’appliquer ici à rechercher, dans tout genre de questions naturelles, la quantité qui produit un maximum ou un minimum : tâche qui semble relever plutôt de la philosophie que des mathématiques.

Puisqu’il existe deux voies pour connaître les effets de la nature — l’une par les causes efficientes, appelée méthode directe, l’autre par les causes finales — le mathématicien emploie l’une et l’autre avec un égal succès. Lorsque les causes efficientes sont trop cachées, tandis que les causes finales échappent moins à notre connaissance, la question se résout généralement par une méthode indirecte ; au contraire, on emploie la méthode directe toutes les fois que l’on peut définir l’effet à partir des causes efficientes.

Il faut surtout veiller à ouvrir l’accès à la solution par les deux voies : ainsi non seulement une solution est confirmée par l’autre, mais encore leur accord nous procure la plus grande satisfaction.

C’est de cette manière que la courbure d’une corde ou d’une chaîne suspendue a été trouvée par deux voies : l’une a priori, à partir des forces de la gravité ; l’autre par la méthode des maxima et minima, puisqu’on comprend que la corde doit prendre une courbure telle que son centre de gravité occupe la position la plus basse.

De même, la courbure des rayons traversant un milieu transparent de densité variable a été déterminée à la fois a priori et à partir du principe suivant : ils doivent parvenir en un temps minimal au point donné.

De nombreux autres exemples semblables ont été proposés par les très célèbres Bernoulli et d’autres savants ; grâce à eux, tant la méthode de résolution a priori que la connaissance des causes efficientes ont fait de grands progrès.


\columnbreak
Quanquam igitur, ob haec tam multa ac praeclara specimina, dubium nullum relinquatur, quin in omnibus lineis curvis, quas Solutio Problematum physico-mathematicorum suppeditat, maximi minimive cujuspiam indoles locum obtineat; tamen saepenumero hoc ipsum maximum vel minimum difficillime perspicitur; etiamsi a priori Solutionem eruere licuisset.

Sic est figura, quam lamina elastica incurvata induit, jam pridem est cognita; tamen quemadmodum ea curva per Methodum maximorum \& minimorum, hoc est, per causas finales, investigari possit, a nemine adhuc est animadversum.

Quamobrem cum Vir Celeberrimus, atque in hoc sublimi naturam scrutandi genere perspicacissimus, Daniel BERNOULLI mihi indicasset se universam vim, quae
in lamina elastica incurvata insit, una quadam formula quam vim potentialem appellet, complecti posse; hancque expressionem in curva Elastica minimam esse oportere; quoniam hoc invento Methodus mea maximorum ac minimorum hoc Libro tradita mirifice illustratur, ejusque usus amplissimus maxime evincitur; hanc occasionem exoptatissimam praetermittere non possum, quin, hanc insignem curvae Elasticae proprietatem a Celeb. Bernoullio observatam publicando, simul Methodi meae usum clarius patefaciam.

Continet enim ista proprietas in se differentia secundi gradus, ita ut ei evolvendae Methodi Problemata isoperimetrica solvendi ante traditae non sufficiant.

\columnbreak
Bien qu’à cause de ces nombreux et remarquables exemples il ne subsiste aucun doute que, dans toutes les lignes courbes fournies par la solution des problèmes physico-mathématiques, intervient un certain maximum ou minimum, il arrive cependant très souvent que ce maximum ou ce minimum soit très difficile à discerner, même si l’on pourrait obtenir la solution a priori.

Ainsi, la forme que prend une lame élastique courbée est connue depuis longtemps ; pourtant, personne n’a encore remarqué comment cette courbe peut être recherchée par la méthode des maxima et minima, c’est-à-dire par les causes finales.

C’est pourquoi, lorsque le très célèbre et très perspicace Daniel Bernoulli m’a indiqué qu’il pouvait exprimer toute la force contenue dans une lame élastique courbée par une certaine formule, qu’il appelle force potentielle, et que cette expression doit être minimale pour la courbe élastique, j’ai vu que cette découverte illustre admirablement la méthode des maxima et minima présentée dans ce livre et en montre très clairement l’immense utilité.

Je ne pouvais donc laisser passer une occasion si favorable sans publier cette propriété remarquable de la courbe élastique observée par Bernoulli, et mettre en même temps plus clairement en lumière l’usage de ma méthode.

En effet, cette propriété renferme des différentielles du second ordre, si bien que les méthodes précédemment exposées pour résoudre les problèmes isopérimétriques ne suffisent pas à son développement.
\end{multicols}

\section*{§ 25.}
\begin{multicols}{2}
Quoniam igitur speciem primam constituimus, si in aequatione generali $dy=\frac{(aa-cc+xx)dx}{\sqrt{(cc-xx)(2aa-cc+xx)}}$ suerit $c=0$, seu $\frac{a}{c}=\infty$, quo casu linea resultat repraesentans statum laminae Elastica naturalem ; ad eandem speciem primam reseramus quoque eos casus, quibus $c$ est quantitas quamminima, ita ut prae $a$ pro evanelsente haberi queat.

Quia ergo $x$ ipsam $c$ superare nequit ; etiam $x$ prae $a$ evanescet, ideoque ista prodibit aequatio $dy=\frac{adx}{\sqrt{2(cc-xx)}}$, cujus integrale est $y=\frac{a}{\sqrt{2}}Asin\frac{x}{c}$, quae est aequatio pro curva Trochoise in infinitum elongata.

Siet autem $AD=\frac{\pi a}{2\sqrt{2}}$, a qua ipsa curvae longitudo infinite parum tantum discrepat, propterea quod angulus DAM est infinite parvus, Sit longitudo laminae ACB$=2f$, ejusque elasticitas absoluta $=Ekk$, ob $f=\frac{\pi a}{2\sqrt{2}}$, erit vis ad hanc curvaturam infinite parvam laminae inducendam requisita finitae magnitudinis \& quidem $=\frac{Ekk}{ff}\frac{\pi \pi}{4}$. Scilicet si extremitates A \& B colligentur silo AB, hoc silum contrahi debebit $vi=\frac{Ekk}{ff}\frac{\pi \pi}{4}$.

\columnbreak
En exploitant l'équation différentielle $\frac{dy}{dx}$=$\frac{(a^2-c^2+x^2)}{\sqrt{(c^2-x^2)*(2a^2-c^2+x^2)}}$ dans le cas où $c=0$, ou $\frac{a}{c}=\infty$ il est possible d'exprimer l'état d'équilibre élastique de la poutre avec $c$ la quantité de ``quamminima'' et $a$ la partie évansecente.\\

$x$ ne peut être supérieur à $c$, dans le cas où $x$ est très petit devant $a$ cela donne l'équation $\frac{dy}{dx}=\frac{a}{\sqrt{2(c^2-x^2)}}$ dont la solution $y=\frac{aA}{\sqrt{2}}sin\frac{x}{c}$ correspond aux poutres infiniment allongées de ``Trochoise''.\\

En posant $AD=\frac{\pi a}{2\sqrt{2}}$, la longueur de la courbe infinie l'angle $DAM$ est infiniment petit. $ACB=2f$ est la longueur de la poutre et son module élastique vaut $E*k^2$ et $f=\frac{\pi a}{2\sqrt{2}}$. Les exigences du calcul infinitésimal appliqué aux poutres s'applique et la force associée vaut $\frac{E*k^2}{f^2}\frac{\pi^2}{4}$. Les extrémités A \& B peuvent être contractées si la force vaut $vi=\frac{E*k^2}{f^2}\frac{\pi^2}{4}$.
\end{multicols}

\end{document}
